\clearpage\pagestyle{otMain}\setlength{\otBlocEspacement}{2.5mm}\fontsize{11}{15}\selectfont\corrigesChapitre{21}{Premiers pas vers le hasard}
\begin{blocLivre}{cor-21-1}\begin{corrige}{1}{Aléatoire ou pas ?}\demarche{Se demander si le résultat est prévisible.}
Le dé et la carte : $\res{\text{aléatoires}}$, le résultat est imprévisible. L’eau qui chauffe : résultat $\res{\text{certain}}$, l’expérience n’est pas aléatoire.
\end{corrige}\end{blocLivre}
\begin{blocLivre}{cor-21-2}\begin{corrige}{2}{Le bon mot}\demarche{Confronter l’évènement à la liste des issues.}
« Obtenir 7 » : $\res{\text{impossible}}$. « De 1 à 6 » : $\res{\text{certain}}$. « Obtenir 2 » : $\res{\text{possible}}$, une issue sur six.
\end{corrige}\end{blocLivre}
\begin{blocLivre}{cor-21-3}\begin{corrige}{3}{Lister les issues}\demarche{Écrire la liste complète sans oubli ni doublon.}
Pièce : $\res{\text{pile, face}}$, 2 issues. Dé : $\res{1,2,3,4,5,6}$, 6 issues. Sac : $\res{\text{rouge, verte, bleue}}$, 3 issues.
\end{corrige}\end{blocLivre}
\begin{blocLivre}{cor-21-4}\begin{corrige}{4}{Une chance sur deux}\demarche{Comparer les issues favorables à la moitié du total.}
Pile : $\res{\text{oui}}$, une issue sur deux. Pair : les issues 2, 4, 6 sur six, soit la moitié : $\res{\text{oui}}$. Carte rouge : 16 rouges sur 32 : $\res{\text{oui}}$ aussi.
\end{corrige}\end{blocLivre}
\begin{blocLivre}{cor-21-5}\begin{corrige}{5}{Sur l’échelle des chances}\demarche{Associer chaque évènement à une zone de l’échelle.}
Janvier après décembre : $\res{\text{certain}}$, à droite. Obtenir 3 : $\res{\text{peu probable}}$, une issue sur six, à gauche du milieu. La boule jaune : $\res{\text{impossible}}$, tout à gauche. Les positions figurent en vert.
\label{fig-cor-21-5}\figureGeo[corrige]{pro-5}
\end{corrige}\end{blocLivre}
\begin{blocLivre}{cor-21-6}\begin{corrige}{6}{Plus ou moins probable}\demarche{Compter les issues qui réalisent chaque évènement.}
« Obtenir 2 » : 1 issue. « Pair » : 3 issues (2, 4, 6). Le pair est $\res{\text{plus probable}}$ : trois fois plus d’issues le réalisent.
\label{fig-cor-21-6}\figureGeo[corrige]{pro-6}
\end{corrige}\end{blocLivre}
\begin{blocLivre}{cor-21-7}\begin{corrige}{7}{L’urne trompeuse}\demarche{Revenir aux boules, les vraies issues du tirage.}
Les issues sont les 4 boules, pas les 2 couleurs : une seule boule donne le rouge, trois donnent le bleu. Le rouge est $\res{\text{peu probable}}$, loin d’une chance sur deux : compter les issues, pas les noms des résultats.
\label{fig-cor-21-7}\figureGeo[corrige]{pro-7}
\end{corrige}\end{blocLivre}
\begin{blocLivre}{cor-21-8}\begin{corrige}{8}{Équiprobable ou non ?}\demarche{Observer si un résultat est favorisé par l’objet.}
Le dé équilibré : $\res{\text{oui}}$, six issues de mêmes chances. La roue aux secteurs inégaux : $\res{\text{non}}$, le grand secteur sort plus souvent. La pièce truquée : $\res{\text{non}}$ plus.
\label{fig-cor-21-8}\figureGeo[corrige]{pro-8}
\end{corrige}\end{blocLivre}
\begin{blocLivre}{cor-21-9}\begin{corrige}{9}{Le sac de la kermesse}\demarche{Comparer à la moitié du nombre total d’issues.}
Cinq issues sur sept donnent le rouge : plus de la moitié, l’évènement est $\res{\text{probable}}$. Le bleu n’a que 2 issues sur 7, moins de la moitié : $\res{\text{peu probable}}$.
\label{fig-cor-21-9}\figureGeo[corrige]{pro-9}
\end{corrige}\end{blocLivre}
\begin{blocLivre}{cor-21-10}\begin{corrige}{10}{Les jetons numérotés}\demarche{Lister les numéros qui réalisent chaque évènement.}
Plus grand que 6 : les issues 7 et 8, soit 2 issues. Plus petit que 4 : les issues 1, 2 et 3, soit 3 issues. Le second évènement est $\res{\text{plus probable}}$.
\end{corrige}\end{blocLivre}
\begin{blocLivre}{cor-21-11}\begin{corrige}{11}{Le jeu de 32 cartes}\demarche{Compter les cartes qui réalisent chaque évènement.}
Les cœurs : 8 cartes. Les rouges : les 8 cœurs et les 8 carreaux, 16 cartes. La carte rouge est $\res{\text{plus probable}}$ : tout cœur est déjà une carte rouge, et les carreaux s’y ajoutent.
\end{corrige}\end{blocLivre}
\begin{blocLivre}{cor-21-12}\begin{corrige}{12}{Gagner ou perdre}\demarche{Prendre les billets comme issues, pas les mots.}
Les issues sont les 500 billets : seuls 10 donnent un lot. Gagner est $\res{\text{peu probable}}$ : deux résultats possibles ne font pas deux résultats de mêmes chances.
\end{corrige}\end{blocLivre}
\begin{blocLivre}{cor-21-13}\begin{corrige}{13}{La roue de la kermesse}\demarche{Compter les secteurs de chaque lot.}
Les secteurs identiques rendent les 8 issues équiprobables. « Perdu » : 4 issues, le résultat le plus probable ; « bonbon » : 3 ; « peluche » : $\res1$ issue sur 8, bien moins qu’une chance sur deux : $\res{\text{non}}$.
\label{fig-cor-21-13}\figureGeo[corrige]{pro-13}
\end{corrige}\end{blocLivre}
\begin{blocLivre}{cor-21-14}\begin{corrige}{14}{La série de six}\demarche{Revenir à la définition d’une expérience aléatoire.}
Les deux sont fausses : le lancer reste une expérience $\res{\text{aléatoire}}$, au résultat imprévisible. Le 6 reste possible, ni certain ni impossible : les lancers passés ne commandent pas le suivant.
\end{corrige}\end{blocLivre}
\begin{blocLivre}{cor-21-15}\begin{corrige}{15}{Classer quatre évènements}\demarche{Compter les issues favorables de chaque évènement.}
Les issues : 10 boules. Rouge : 2, verte : 3, noire : 5, et le dernier évènement est réalisé par les 10 : il est $\res{\text{certain}}$. Classement : $\res{\text{rouge}<\text{verte}<\text{noire}<\text{certain}}$.
\label{fig-cor-21-15}\figureGeo[corrige]{pro-15}
\end{corrige}\end{blocLivre}
